% conclusion.tex
% Language-audit pass, 2026-07-14 (see project_language_audit_2026-07-14
% memory). Reframed 2026-07-14 so the multi-asset variance correction
% leads and the HMM-WJ result supports the marginal component. The
% conclusion ends on the missing cross-asset covariance limitation.
% 2026-08-31: merged to two paragraphs (contributions and results, then
% the cross-asset covariance limitation) to match the shortened
% Discussion and the journal's "short Conclusions section" guidance.

In this study, we developed a variance-corrected single-index
composition method that reuses full-return generators in a factor model
without double counting their location or variance. Centering first
prevented a full-return draw from shifting the calibrated intercept;
the subsequent scale correction removed the variance inflation caused
by adding the market factor. Across $423$ non-market assets,
it recovered the calibrated factor loading to a low median absolute
error and produced heavy-tailed marginals. On $416$ complete held-out
asset histories from $2025$, it improved the mean cross-ticker KS
pass rate over centered naive composition and held the median
variance ratio near one, while naive composition remained inflated.
Its one-day left-tail $99\%$ Value-at-Risk exceedance rate matched
that of JumpHMM-on-residuals to the reported precision but stayed
above the nominal rate. Those residual models,
fitted directly to the OLS residual series, required a separate fit
for each asset and produced a similar marginal pass rate.
Solving the composition problem first required a per-asset generator
that retained heavy tails and temporal structure, for which we used a
hidden Markov marginal model with Laplace-quantile states,
Student-$t$ emissions, and direct transition counting. On SPY, its jump-duration variant reduced
absolute-return autocorrelation error relative to the no-jump variant
while keeping high KS and AD pass rates; GARCH reached a lower
autocorrelation error but failed the distributional tests. The
single-asset result therefore showed a measured tradeoff between
distributional fit and volatility clustering, not uniform superiority.
Taken together, these results separate two tasks that are easily
conflated when synthesizing multi-asset market data: the marginal
generator governs each asset's distributional and temporal behavior,
whereas the variance correction prevents the common market factor from
double counting the generator's total variance.

The current multi-asset implementation does not model cross-asset
residual covariance. Independent residuals omit the sector and style
dependence that links real assets, and a rebalanced portfolio built on
them earns a diversification return that a real portfolio would not.
A full joint residual model is needed before using these paths for
portfolio-level return forecasts. The variance correction is a
necessary compositional constraint rather than a complete dependence
model. Portfolio-level synthetic market data would require combining
it with joint residual dynamics.
